\documentclass[../main.tex]{subfiles}
\begin{document}
%==========================================TIÊU ĐỀ TẬP========================================
\begin{table}[h]
    \begin{adjustwidth}{-1cm}{}
        \begin{tabular}{>{\centering\arraybackslash}p{6cm}|>{\raggedright\arraybackslash}p{12cm}}
            \multirow{5}{6cm}
            % Cột bên trái: ÉPISODE và số tập
            {\\[-40pt]
                \centering
                {\fontsize{20pt}{20pt}\selectfont{\outlinetext[1]{eptype}{white}{\flaregothic ÉPISODE}}\\[12pt]
                    {\fontsize{72pt}{72pt}\selectfont{\outlinetext[1]{epnum}{white}{\burbank 166}}}\\[6pt]       % use for normal/special episodes
                    % {\fontsize{36pt}{36pt}\selectfont{\outlinetext[1]{epnum}{white}{\burbank SÁU}}}\\[6pt]     % use for custom episodes
                    {\fontsize{20pt}{20pt}\selectfont{\outlinetext[1]{epnum}{white}{\cafeteria (13-7)}}}\color{black}}
                \\[18pt]
                {\fontsize{14pt}{14pt}\selectfont{\outlinetext[1]{epattr}{white}{\flaregothic BIENVENUE, }}{\flaregothic\color{Luilui} Lui}\fontsize{14pt}{14pt}\selectfont{\outlinetext[1]{epattr}{white}{\flaregothic \ !}}}\\
                % {\fontsize{18pt}{18pt}\selectfont{\outlinetext[1]{epattr}{white}{\flaregothic [II]}}}\\
                \color{black}
            }
             &
            % Dòng thứ nhất: tên tập tiếng Nhật
            {\fotskip\fontsize{18pt}{18pt}\selectfont \ruby{余}{よ}\ruby{裕}{ゆう}ぶってるけど\ruby{全}{ぜん}\ruby{然}{ぜん}お\ruby{金}{かね}がない！}             \\[6pt]
            % Dòng thứ hai: tên tập tiếng Anh
             & {\cafeteria\fontsize{18pt}{18pt}\selectfont Going for Broke}                                                               \\[4pt]
            % Dòng thứ ba: tên tập tiếng Việt
             & {\vnmsans\fontsize{18pt}{18pt}\selectfont Chiến thuật chi tiêu}                                                            \\[8pt]
            % Dòng thứ tư: Nhân vật xuất hiện
             & {\caviar{\fontsize{14pt}{14pt}\selectfont Nhân vật: \emp{kuon2} \emp{kaiho2} \emp{game} \emp{aki} \emp{subaru} \emp{lui}}}
            %  \\
            % Dòng thứ năm (nếu có): Nhân vật chỉ xuất hiện trong Omake
            %  & {\abv Truyện phụ:} \emp{laplus}
            \\[8pt]
            % Dòng cuối cùng: Thời gian diễn ra của tập (thường sẽ là ngày phát hành)
             & {\sen\fontsize{16pt}{16pt}\selectfont Mốc thời gian: 17/07/2022}                                                           \\[18pt]
        \end{tabular}\\[8pt]
    \end{adjustwidth}
\end{table}
%=========================================TRUYỆN CHÍNH========================================
\color{black}

\txtbx[2]{
    {\color{vol} \uline{Tóm tắt:}} Trùng vào ngày sinh nhật lần thứ 26 của cặp song sinh Kuon và Kaiho, Takane Lui tới xin tư vấn tháo gỡ bài toán tài chính chạm vạch báo động đỏ của holoX. Sau khi chiêu trò kêu gọi tài trợ bị Subaru đập tan, Lui chuyển sang màn ``lùa gà'' bán lá nhựa baran cho cô nàng ngây thơ Aki Rosenthal khiến văn phòng náo loạn. Sau khi cuộc hỗn loạn được dập tắt Lui bất ngờ tặng bánh sinh nhật và vô tình để lộ manh mối bí mật về ``bộ ba nhà Chó'', trong khi La+ Darknesss gửi tin nhắn thách đấu nhưng lại chứa đầy sự dễ thương.
}

\pov{Hikawa Kuon}

Hôm nay đánh dấu sinh nhật lần thứ hai mươi sáu của tôi và Kaiho.

Ngày này lại trùng đúng vào tuần thứ tư trong chuỗi năm tuần giới thiệu các thành viên của holoX. Thật trớ trêu thay, ngay cả vào ngày sinh nhật của mình, tôi vẫn phải đi làm như mọi ngày khác. Dù vậy, tôi đã quá quen với điều này rồi.

Được rồi, để chuyện sinh nhật sang một bên đã, bây giờ đến lúc tôi giới thiệu thành viên tiếp theo của hội kín đầy tiếng tăm này.

Người đứng thứ hai tổ chức bí mật holoX, người luôn mang trong mình vẻ ngoài trưởng thành, thanh lịch và đầy sức hút. Nếu Koyori là nhà khoa học luôn tò mò với những phát minh kỳ lạ, thì Lui chính là người lãnh đạo tài năng, người giữ cho toàn bộ tổ chức đi đúng hướng và kéo những thành viên hơi lập dị của mình về đúng khuôn khổ.

Dù thường xuất hiện với hình ảnh một tổng trưởng nghiêm khắc, giữ vững nguyên tắc, nhưng thực tế cô lại là người rất dễ gần và luôn quan tâm đến mọi người xung quanh. Nhờ khả năng lãnh đạo, kỹ năng giao tiếp khéo léo và sự nhạy bén trong công việc, cô luôn nỗ lực hoàn thành vai trò của mình một cách tốt nhất, ngay cả khi đối mặt với những tình huống bất ngờ không lường trước.

Nhưng đừng bị vẻ ngoài chuyên nghiệp đó đánh lừa: tính cách của Lui còn có rất nhiều điểm thú vị khác. Khi ở bên những người thân thiết như Koyori, Chloe hay Iroha, cô trở thành người chị cả đáng tin cậy, luôn chăm sóc và chiều chuộng các thành viên trong tổ chức. Tuy nhiên, đôi khi cô cũng khá vụng về, dễ bị mọi người trêu chọc và đặc biệt là rất thích những thứ dễ thương. Một tổng trưởng tưởng chừng luôn nghiêm nghị, nhưng nhiều lúc lại chẳng khác gì một bà mẹ đang phải chạy đua trông nom đám con nít nghịch ngợm.

Chưa hết, cô còn là một chuyên gia trong lĩnh vực ca hát và chơi guitar.

Và dù là vị trí phó holoX, mục tiêu của cô không chỉ dừng lại ở việc giữ cho tổ chức hoạt động ổn định. cô mong rằng một ngày nào đó sẽ đưa holoX trở thành một cái tên lớn, đồng thời khẳng định năng lực của bản thân trên sân khấu với tư cách là một thần tượng thực thụ.

Một nữ tổng trưởng với khát vọng chinh phục đỉnh cao.

\model{25}{l}{7.5cm}{lui}

\dots Hay có lẽ, đó chỉ là một ``chị đại'' bất đắc dĩ, người đang phải đau đầu quản lý cả một tổ chức mà toàn những kẻ chẳng bao giờ chịu nghe lời!

\chrint

Vâng, đó chính là Takane Lui (\ruby{鷹}{たか}\ruby{嶺}{ね}ルイ, \ruby{Ta-ca}{Ưng}-\ruby{nê}{Lĩnh}\ Lư-i), Phó Tổng trưởng của Tổ chức Bí mật \textit{holoX}.

Với vai trò là đầu mối liên lạc chính và người phụ trách các vấn đề đối ngoại của holoX, cô thực chất là người lãnh đạo thực tế của nhóm, đảm nhận những công việc mà người sáng lập không thể giải quyết được. Dù thoạt nhìn có vẻ lạnh lùng và khó gần, nhưng thực tế Lui lại là một người rất có lòng trắc ẩn, luôn quan tâm sâu sắc đến các cấp dưới của mình. Tiếc thay, cô cũng khá dễ mắc sai lầm đúng vào những khoảnh khắc quan trọng nhất.

Lui có thể được mô tả là kiểu người ``chị cả'', với tính cách nói chung dịu dàng nhưng đôi khi cũng rất năng động và vui vẻ. Nhờ bản tính tương đối bình tĩnh và trưởng thành, cô thường phải gánh vác nhiệm vụ kiềm chế những hành động lập dị, hỗn loạn của các thành viên cùng thế hệ, đặc biệt là La+ và Chloe. Giống như Ninomae Ina'nis \footnote{Sẽ xuất hiện chính thức ở quyển UMISEA và quyển Mười Ba.} và Ouro Kronii\footnote{Sẽ xuất hiện ở quyển Mười Bảy.}, cô cũng rất thích những câu đùa ``bố già'' ngớ ngẩn.

Cô cũng rất dễ bị sợ hãi và hoảng loạn trước bất kỳ điều gì bất ngờ; một ví dụ nổi tiếng là khi cô phát trực tiếp chơi game Outlast vào những ngày đầu hoạt động, với thử thách tự đặt ra là sẽ kết thúc buổi stream ngay lập tức nếu mình la hét. Kết quả là cô phải chơi thêm hai lượt nữa, sau khi lượt đầu tiên thất bại chỉ vì vô tình làm đổ cốc nước.

Trong số các thành viên cùng nhóm, Lui là người tích cực nhất trong việc giao tiếp với các thành viên \hll\ ở nước ngoài cũng như người hâm mộ quốc tế.

Về ngoại hình, cô có mái tóc ngắn màu hồng cùng những chiếc cánh nhỏ màu hồng trên đầu, đôi mắt xanh (có thể chuyển thành màu cam với đồng hình ngắm bắn), và có vòng một đầy đặn. Cô còn có đôi chân dài, hông rộng với khoảng trống giữa đùi khá rõ nét, tạo ra một vóc dáng đặc biệt mà nhiều người hâm mộ ví như một chiếc nĩa điều chỉnh âm thanh. Kiểu tóc mặc định của cô là tóc ngắn được chia ngôi, giữ bằng băng đô đen, cùng mái tóc hai lớp rơi xuống giữa hai mắt.

Trang phục của cô bao gồm: một chiếc vòng cổ đen, áo sơ mi trắng dài tay có khóa kéo ở ngực, cà vạt đỏ với khăn xếp trắng bên dưới, dây đeo ngực dưới ngực, kẹp tay áo, găng tay đen không ngón, quần soóc da đen cao eo, vớ đen, dây đeo đùi bên phải có treo một chiếc roi, và giày cao gót đen quai mảnh. Lui còn có một chiếc áo khoác viền lông vũ đỏ tùy chọn thường khoác trên vai, cùng cặp kính râm tròng cam.

\chrint

Cô nàng là một trong những thành viên mà tôi là đánh giá là ``ngoan ngoãn'' nhất hội. Còn tại sao tôi lại nói thế? Để tôi kể cho các bạn nghe.

\il

Một buổi tối tĩnh mịch nọ tại văn phòng chính của hololive.

Khi phần lớn nhân viên đã lục đục ra về, bên trong phòng làm việc vẫn còn le lói ánh đèn bàn. Takane Lui - người được mệnh danh là ``chị đại đỡ đầu'' kiêm quản lý sổ sách của tổ chức bí mật holoX - đang cặm cụi ngồi thống kê lại toàn bộ thu chi trong hội. Ngồi cạnh em ấy lúc này là anh em song sinh nhà Hikawa: tôi đang chống cằm xem xét báo cáo, còn Kaiho ngồi vắt chân bên cạnh, tay cầm bút gõ nhịp cành cạch xuống mặt bàn với ánh mắt sắc lẹm đầy soi xét.

\img{13-7a}

Đúng vậy, hội kín của em ấy cũng phải vận hành thu chi ngân sách rành mạch chẳng khác gì cách mà tôi và A-san thường làm cho cả tập đoàn này. Thế nhưng, nếu những bản báo cáo tài chính của chúng tôi luôn tập trung kiểm kê các khoản đầu tư chính đáng như chi phí đào tạo thanh nhạc, mua sắm trang thiết bị, tu sửa cơ sở vật chất, hay tổng doanh thu từ hội viên và vật phẩm lưu niệm\dots

Lui thở dài, đẩy tập tài liệu sang phía chúng tôi: ``\dots Đây là toàn bộ biểu đồ những gì mà holoX chúng em đã chi tiêu trong suốt tháng vừa qua.''

Tôi nheo mắt lướt qua từng con số: ``Hừm\dots''

Kaiho liếc mắt nhìn vào bảng kê khai, nhếch mép cười khẩy: ``Nhìn qua cơ cấu phân bổ ngân sách này là em thấy có mùi bất ổn rồi đấy.''

\dots thì bản báo cáo của Hội kín kia, bên cạnh những khoản thu chi hợp lý hiếm hoi, lại chình ình hàng loạt danh mục tiêu xài cực kỳ\dots\ lặt vặt và ba chấm. Cụ thể như sau:

\img{13-7b}

\begin{itemize}
    \item 29\% là chi phí mua gà rán;
    \item 12\% chi phí thức ăn vặt;
    \item 17\% ngân sách phục vụ nghiên cứu;
    \item 8\% chi phí rèn luyện thể lực;
    \item 10\% phí phục vụ nhu cầu giải trí;
    \item 5\% chi phí mua cà tím;
    \item 7\% trích lập quỹ dự phòng;
    \item 12\% là các khoản chi phí phát sinh nhỏ lẻ khác.
\end{itemize}

Tóm lại một câu ngắn gọn, gần một phần ba tổng ngân quỹ của cả hội đã bị đốt sạch vào những thứ vô cùng ngớ ngẩn. Khoản gà rán thì tôi còn có thể tạm hiểu cho cái dạ dày không đáy của ai đó, nhưng đến cả\dots\ cà tím ư? Tại sao lại tốn tận 5\% ngân sách chỉ để mua cà tím cơ chứ?

Tuy nhiên, điều khiến tôi bất ngờ nhất là tỉ lệ nợ nần của họ vẫn giữ ở mức 0\%. Dĩ nhiên chúng ta sẽ tạm bỏ qua khoản cước phí mạng internet cũng bằng 0\%, bởi lẽ cả hội đang ngang nhiên xài chùa đường truyền Wi-Fi của công ty. Dù sao đi nữa, căn cứ theo biểu đồ biến động chi tiêu trong vòng 12 tháng gần nhất, tình hình tài chính của holoX đang rơi thẳng vào ngưỡng báo động đỏ suốt hai tháng liên tiếp.

Lui đưa tay ôm trán ngao ngán: ``Lại thêm một tháng nữa chạm vạch báo động đỏ rồi à?''

Tôi tặc lưỡi, chỉ tay vào biểu đồ tròn: ``Anh thấy gần một nửa số tiền trong này mấy đứa toàn ném qua cửa sổ vào mấy thứ đẩu đâu ấy\dots''

Kaiho khoanh tay lạnh lùng bổ sung: ``Nếu là sổ sách bên bộ phận của bên nhà chị mà lòi ra mấy cái mục `gà rán' với `cà tím' chiếm sóng thế này thì chắc chị cho cả lũ nhịn cơm tối trừ nợ từ lâu rồi.''

Lui cười trừ bất lực: ``Em cũng chẳng dám trông mong kỳ vọng gì nhiều ở ý thức tiết kiệm của mọi người trong hội lắm đâu. Mỗi người đều có một sở thích chi tiêu riêng khó bỏ mà.''

Từ bên trong chiếc đồng hồ đeo tay của tôi, trợ lý ảo Game lười biếng cất giọng phân tích: ``Để tôi thử đoán xem nào. Khoản nghiên cứu chắc chắn là của Hakui-san, chi phí luyện tập thuộc về Kazama-san, còn mục thức ăn và gà rán khổng lồ kia là của Sakamata-san với sếp của cô, hoàn toàn chính xác chứ?''

Lui gật đầu cái rụp: ``Chính xác đến từng milimet luôn ạ. Bản thân em thì hầu như chẳng mấy khi chi tiêu hoang phí cho cá nhân.''

Tôi nhướn mày sâu hơn, chắp hai tay lại suy luận: ``Vậy nói tóm lại\dots\ toàn bộ phần chi tiêu lặt vặt quái đản còn lại chiếm gần nửa ngân sách kia đều xuất phát từ vị Sếp tổng của các em, đúng chứ?''

Nàng Chim ưng khẽ gật đầu xác nhận kèm theo một tiếng thở dài sườn sượt. Tôi và Kaiho nhìn nhau đổ mồ hôi hột.

Tôi buột miệng thốt lên: ``Thật không thể tưởng tượng nổi một người mang danh nghĩa sáng lập tổ chức lại có thể tiêu xài hoang phí vô độ đến mức độ này đấy.''

Kaiho nhăn mặt khinh bỉ: ``Cái mớ hỗn độn đó mà cũng gọi là chi tiêu điều hành à? Đúng là một đứa trẻ to xác phá gia chi tử.''

Game tiếp tục đặt câu hỏi thực tế: ``Vậy hiện tại cô đã vạch ra được giải pháp khắc phục nào khả thi chưa?''

Tôi gợi ý: ``Ví dụ như việc siết chặt kỷ luật, cân đối lại dòng tiền thu chi chẳng hạn? Hay là ban hành chính sách thắt lưng buộc bụng tiết kiệm triệt để?''

Lui khoanh tay trước ngực, nhắm mắt trầm ngâm suy tính về bài toán tài chính nan giải này. Sau một hồi vắt óc động não căng thẳng, đôi mắt em ấy chợt lóe sáng: ``Bọn em sẽ cần phải chủ động tìm kiếm thêm những nguồn thu nhập mới!'' Đó xem ra cũng là một hướng đi có tính đột phá.

Game hỏi vặn: ``Thế các cô định kiếm nguồn thu đó từ đâu ra?''

Lui nghiêng đầu đẩy bóng sang phía chúng tôi: ``Hừm\dots\ Hai anh chị thử vắt óc nghĩ cách giúp em xem nào.''

Tôi vội vàng xua tay từ chối: ``Đừng có mà đùn đẩy hết trách nhiệm sang cho anh. Anh có phải là thành viên biên chế thuộc hội kín của mấy đứa đâu.''

Kaiho cũng lạnh lùng phẩy tay: ``Tôi bận tối mắt với đám báo cáo của Holostars rồi, tự làm tự chịu đi.''

Sau một hồi tự mình tiếp tục suy tư cao độ, cuối cùng Lui cũng đã vạch ra được một kế hoạch tác chiến hoàn chỉnh (mà dĩ nhiên là tôi có thể đọc thấu suy nghĩ trong đầu em ấy rồi, nhớ chứ?). Đôi mắt của nàng Chim ưng bỗng chốc biến đổi sắc tố, chuyển sang màu vàng rực rỡ với biểu tượng chữ thập màu đỏ hồng sắc lẹm, khác hẳn với màu xanh nước biển dịu dàng thường ngày. Chắc hẳn em ấy đã tìm ra được chìa khóa giải quyết bài toán hóc búa này rồi đây. Nói trước luôn, mấy cái ý tưởng nảy ra trong đầu em ấy lúc này cũng chẳng hay ho tử tế gì cho cam đâu. Nhưng thôi, tôi quyết định giữ im lặng để xem kịch hay.

\img{13-7c}

Lui dõng dạc tuyên bố: ``Quyết định vậy đi! Bắt đầu triển khai kế hoạch thôi!''

Ngay lập tức, một màn hình trình chiếu cỡ lớn đột ngột hiện ra lơ lửng giữa không trung, bắt đầu phát đoạn video giới thiệu hoành tráng về tổ chức bí mật holoX. Mở màn lần lượt là chân dung của Laplus, Lui, Koyori--

*TÁCH!*

Toàn bộ hệ thống đèn điện trong văn phòng bỗng chốc bật sáng choang. Tôi giật mình quay đầu lại nhìn về phía cửa ra vào và bắt gặp Subaru đang đứng sừng sững ngay bên cạnh công tắc đèn với vẻ mặt ngơ ngác.

\img{13-7d}

Game cất giọng chào: ``Ồ, chào Oozora-san.''

Subaru chớp mắt ngạc nhiên: ``Ê-tồ\dots\ Kuon-senpai, Kaiho-san, rồi cả Lui-chan nữa, mọi người đang tụ tập làm cái trò gì trong bóng tối thế này vậy?''

Tôi thản nhiên đáp lại: ``Anh cũng đang định mở miệng hỏi ngược lại em câu đó đấy.''

Kaiho khoanh tay tựa lưng vào ghế, nhếch mép: ``Lại đi tuần tra đêm khuya thay bảo vệ tòa nhà đấy à, Oozora?''

Do sự xuất hiện quá đỗi bất ngờ của nàng Vịt, cô nàng tóc nâu Lui lập tức giật bắn mình, vội vàng thu lại ánh mắt sắc lẹm về trạng thái bình thường và luống cuống chữa cháy: ``À, em\dots\ em xin lỗi vì đã làm phiền ạ!'' Nói đoạn, em ấy lật đật ôm vội chiếc laptop của mình né sang một góc khác đầy lén lút.

Tôi chỉ biết đổ mồ hôi hột trước hành động mờ ám có phần giấu đầu lòi đuôi của em ấy. Thôi được rồi, đằng nào cũng đã trót dở dang, để tôi tiếp tục nốt màn giới thiệu cho Hội kín của em ấy vậy. *E hèm!*

Tôi hắng giọng dõng dạc cất lời: ``Và mở đầu cho đội hình tinh nhuệ này chính là Người sáng lập tối cao của Hội: La+ Darkness--''

``\textbf{Gượm cái mồm lại đã nào!!}'' lập tức, nàng Tomboy lao vụt tới chen ngang lời bình luận của tôi, tiện tay vung chưởng đập vỡ tan tành cái màn hình trình chiếu ảo ảnh kia. Cô nàng chống hông quát lớn: ``\textbf{Mấy cái thứ màu mè này rốt cuộc là cái quái gì đây hả!?}''

\img{13-7e}

Game bình thản châm chọc: ``Thật trùng hợp là tôi cũng đang định đặt câu hỏi y hệt như vậy cho Takane-san đấy.''

Lui gãi đầu, nở một nụ cười trừ gượng gạo: ``Hì hì\dots\ Em nghĩ là đã đến lúc thích hợp để phát sóng phần mở đầu giới thiệu nhân vật rồi đấy chứ ạ.''

Subaru trợn mắt phản đối: ``Mở đầu cái con khỉ khô gì chứ! Đừng có mà tự tiện phát sóng mấy cái thứ ngớ ngẩn sặc mùi chuunibyou như thế nữa đi!''

Tôi tặc lưỡi tiếc nuối: ``Chán thật sự luôn chứ\dots\ Phần mở đầu đang vào đoạn cao trào hay ho thế mà em lại nỡ phá hỏng.''

Kaiho liếc nhìn tôi với ánh mắt khinh bỉ: ``Anh bớt hùa theo mấy cái trò trẻ con nhảm nhí đấy giùm em đi, Aniki.''

Subaru quay phắt sang gào lên với tôi: ``Anh đang nói nhảm nhí cái gì thế hả, Kuon-senpai!?''

Bất chợt, bầu không khí xung quanh Lui trầm hẳn xuống, cô nàng tỏ rõ thái độ không vui ra mặt. Nét mặt em ấy trở nên nghiêm nghị đến đáng sợ, cất giọng lạnh lùng đầy bất mãn: ``Vậy à? Nếu chị đã chê bai thẳng thừng như thế, vậy thì phiền Subaru-senpai đứng ra làm phần mở đầu quảng bá thay cho bọn em đi nhé?''

Cả nàng Vịt lẫn trợ lý ảo của tôi đều đồng thanh hét lên kinh ngạc: ``Cái thứ logic quái quỷ gì thế này!?''

Ngay khoảnh khắc đó, hai dòng chữ quảng cáo to đùng bỗng hiện lên lù lù ngay trước mặt họ:

\begin{center}
    \uline{``Công ty khai thác Haachama\\
        Xin trân trọng giới thiệu''}
\end{center}

\img{13-7f}

Tôi thừa nhận đó cũng là một phương thức kiếm tiền tài trợ khá hữu dụng trong thời buổi kinh tế khó khăn này, cơ mà tôi nghĩ đây tuyệt đối không phải là thời điểm thích hợp để em ấy giở chiêu trò đó ra đâu. Thậm chí, Subaru còn nhanh tay vung chưởng đập nát bét hai dòng chữ quảng cáo đang bay lơ lửng xuống đất, khiến Lui cũng chẳng kịp trở tay đỡ đòn.

Nàng Vịt gào lên đầy bức xúc: ``Với cả, làm ơn dẹp ngay mấy cái thứ chữ nghĩa nhằng nhịt này ra khỏi mặt chị ngay lập tức!''

Lui mếu máo tiếc nuối: ``Không chịu đâu! Đấy là suất quảng cáo kiếm cơm của bọn em mà!''

Tôi lắc đầu thở dài: ``Giờ không phải là lúc thích hợp để chạy quảng cáo thương mại cho Hội đâu em ơi\dots''

Kaiho đứng bên cạnh khoanh tay cười nhạt: ``Chưa kịp thu được đồng vốn nào đã bị đập nát bét thế này thì đúng là lỗ nặng rồi, Lui ạ.''

Nàng tóc đen Subaru quay ngoắt đầu về phía anh em chúng tôi, trố mắt hỏi dồn: ``Mọi chuyện rốt cuộc là sao đây, Kuon-senpai, Kaiho-senpai? Chuyện quái quỷ gì đang diễn ra ở cái văn phòng này vậy!?''

Tôi chỉ biết thở dài sườn sượt, hất hàm đáp lại: ``Đi mà hỏi cái cô nàng đang cố sống cố chết kéo doanh thu quảng cáo ở đằng kia kìa.''

Nàng Chim ưng khẽ thở dài sườn sượt, cuối cùng cũng chịu mở lòng bộc bạch: ``Chẳng là\dots\ tình hình tài chính hiện tại của holoX đang rơi vào cảnh khó khăn lắm ạ\dots\ Nói một cách ngắn gọn thì đúng là cực kỳ gian nan luôn.''

Đó chính xác là những gì mà anh em chúng tôi vừa tận mắt chứng kiến trên bản biểu đồ thu chi chi chít báo động đỏ mà em ấy vừa đưa xem.

Tiền nong suy cho cùng thực sự là một vấn đề nhức nhối khiến cho người ta đau đầu đến bạc tóc. Gia đình nhà Hikawa chúng tôi ngày trước cũng từng vài ba lần rơi vào cảnh túng thiếu như vậy, và cứ mỗi lần khủng hoảng tài chính ập đến là cả nhà lại phải gồng mình thắt chặt chi tiêu từng đồng từng cắc. Khó khăn và áp lực lắm chứ chẳng đùa.

Thế nhưng, nỗi lo cơm áo gạo tiền ấy thực sự vẫn còn đỡ đau đầu hơn nhiều so với việc phải cố gắng vắt óc giải mã tại sao Lui lại có thể tận dụng chính cái sự khốn khổ của bản thân chỉ để\dots\ chơi chữ tấu hài. Tôi và Kaiho thực sự bị dị ứng và ám ảnh với cái kiểu đùa nhạt nhẽo đó nhất trên đời.

Lui nghiêng đầu lẩm bẩm: ``Gian nan\dots?''

Trợ lý Game đáp gọn: ``Thì chính miệng cô vừa mới thốt ra từ đó đấy thôi.''

Và thế là bánh răng vô tri trong đầu em ấy bắt đầu xoay chuyển\dots

Lui sáng mắt lên: ``Hmm? Tình hình đang gian nan? Mà vấn đề lại nằm ở tiền\dots\ tiền nong\dots'' rồi như vừa nghĩ ra ý tưởng nào đó, một bóng đèn điện liền bật sáng trưng trên đầu: ``À, vậy ghép lại thì sẽ thành\dots\ nan nong!''

Tôi há hốc mồm, trố mắt ngơ ngác: ``\dots Hả?''

Kaiho đứng bên cạnh nhíu mày sâu hoắm, gân xanh giật giật liên hồi trên trán: ``Lại là cái trò chơi chữ nhạt như nước ốc này nữa hả trời!? Lạy em đấy Lui! Hết cái vụ `Sushi' bỏ chữ S thành `Ushi' rồi biến thành trâu bò nhảy disco náo loạn cả phố ở tập 153\footnote{{\flaregothic ÉPISODE 153}, quyển số Mười Hai.}, rồi đến cái màn tra tấn vần `-zai` từ `oan sai', `vạn tuế' cho tới `nước rửa chén' ép Aru-ru biến hình liên tục bên căn mancave\footnote{{\humanist FOLGE 16}, quyển số Hai thuộc \textit{stars*collection}.}\dots\ Giờ em lại còn lôi cả cái bánh naan Nam Á vào để chơi chữ `gian nan' với `nan giải' nữa! Bộ mấy người ở cái công ty này bị nhiễm virus nghiện chơi chữ giai đoạn cuối hết rồi đấy à!?''

Game từ cổ tay tôi cũng thở dài ngao ngán, buông một lời đả kích: ``Tôi đề xuất Kaiho-user dùng vũ lực dập tắt ngay mầm mống chơi chữ này trước khi nó kịp lây lan sang Oozora-san. Chỉ số nhạt nhẽo trong phát ngôn vừa rồi của Takane-san đã vượt quá ngưỡng chịu đựng tối đa của bộ xử lý lượng tử rồi.''

Tôi rùng mình xoa thái dương: ``Công nhận\dots\ Nhắc lại hai cái vụ kia là anh vẫn còn thấy đau hết cả tiền đình đây này\dots''

Ngay lập tức, nàng Vịt Subaru không nói nhiều, vung tay ném thẳng một chiếc bánh cỡ bự giáng thẳng xuống đỉnh đầu của cô nàng ``thánh chơi chữ'', đập bẹp dí em ấy xuống sàn nhà. ``\textbf{Đừng có mà giở cái thói chơi chữ nhảm nhí đó ra nữa!}''

Tôi ngơ ngác chỉ tay: ``Ờm\dots\ Cái vật thể lạ khổng lồ gì đang nằm trên đầu em ấy đây?''

Lui lồm cồm bò dậy từ dưới đống bánh, chẳng những không giận mà còn cười tít mắt: ``Chị lại hào phóng tặng cho em một chiếc bánh naan\footnote{Một loại bánh mì dẹt truyền thống nướng lò của vùng Nam Á, chủ yếu là ở Ấn Độ.} cơ à\dots\ Chơi chữ hay đấy chứ!''

\img{13-7g}

Game thở dài sườn sượt qua cổ tay: ``Rơi vào hoàn cảnh bi đát gian nan thế này rồi mà cô vẫn còn tâm trí đùa cợt được thì tôi cũng chịu thua\dots''

Dứt lời, nàng Chim ưng thản nhiên ôm chầm lấy chiếc bánh mì khổng lồ kia rồi đưa lên miệng cắn xé nhai một cách ngon lành. Nhìn cảnh tượng đó, tôi càng cảm thấy mù mịt, hoàn toàn không hiểu nổi ý đồ thực sự ẩn sau hành động của em ấy này là gì.

Subaru ôm đầu bất lực: ``Thật đúng là một sinh vật quái gở kỳ quặc hết thuốc chữa.''

Tôi lắc đầu ngao ngán: ``Anh thậm chí còn phải tự vứt não đi thì mới mong tiêu hóa nổi những gì em ấy vừa nói. Với cả, hai đứa bây kiếm đâu ra cái chiếc bánh to chà bá lửa như cái mâm đồng vậy?''

Kaiho khoanh tay châm chọc: ``Chắc là dùng cả 12\% ngân sách chi phí phát sinh để đặt lò bánh mì làm riêng rồi.''

Cô nàng Chim ưng hào hứng: ``Senpai có biết không ạ? Loại bánh này ở tận bên khu vực Nam Á người ta cũng chẳng bao giờ làm to đến mức ngoại cỡ thế này đâu.''

Cô nàng tóc nâu đứng phắt dậy, hai tay vẫn ôm khư khư chiếc bánh khổng lồ, ánh mắt bỗng chốc ánh lên vẻ cảm thông sâu sắc dành cho\dots\ một ông chủ nhà hàng cà-ri trời ơi đất hỡi nào đó, khi họ đã phải còng lưng vất vả nhào nặn ra một chiếc bánh siêu to khổng lồ như vậy.

``Họ chắc hẳn đã phải cất công xây hẳn một cái lò nướng to bự chảng như thế này để phục vụ thực khách đấy\dots'' trong lúc em ấy còn đang say sưa thả hồn trải hết cảm xúc sướt mướt vào món ăn truyền thống của một trong những cái nôi văn hóa nhân loại, cả Game, tôi, Kaiho lẫn Subaru đều đồng loạt vội vàng lao tới chặn họng em ấy lại.

Tôi xua tay lia lịa: ``Anh biết tỏng là em đang muốn bày tỏ tình cảm dạt dào của mình vào chiếc bánh này rồi. Gớm nào là gian nan, rồi nan giải, rồi\dots\ đủ các loại `nan' trên đời nữa, nhưng mà làm ơn\dots''

Subaru cau có quát lớn: ``Em đừng có mở miệng nói thêm mấy cái thứ tào lao bí đao như vậy nữa được không hả!?''

Game nhắc nhở: ``Tập trung quay trở lại vấn đề chuyên môn tài chính đi giùm cái.''

Kaiho bẻ khớp ngón tay *rắc rắc*, lườm nguýt: ``Mở miệng chơi chữ thêm một câu nữa là chị nhét nguyên cái bánh này vào họng em đấy, Takane.''

Lui chớp mắt ngây thơ: ``À, ý các anh chị là dẹp bỏ mấy cái thứ nhan nhản buồn chán như vậy đi đúng không ạ?''

Subaru hoảng loạn lao tới cố gắng lấy tay bịt chặt miệng Lui lại: ``Đã bảo là đừng có chơi chữ nữa mà! Ngậm cái miệng lại ngay lập tức đi!''

Tôi cạn sạch ngôn từ: ``Anh thực sự không tài nào hiểu nổi cái cấu trúc ngôn ngữ trong đầu em đang vận hành kiểu gì luôn đấy\dots''

Gần như ngay sau khoảnh khắc hỗn loạn đó, Aki liền xông thẳng vào văn phòng bằng một cú đạp chân trời giáng khiến chiếc cửa gỗ gãy rời cả bản lề, bay thẳng vào trong. ``Em vừa mới gọi chị đến trợ giúp đúng không hả, Subaru-chan?'' em ấy hớn hở cất giọng, giơ ngón tay cái lên đầy tự tin.

\img{13-7h}

Tôi nhìn cái cánh cửa tội nghiệp đang nằm chỏng chơ dưới sàn mà chỉ muốn lao tới nhéo tai em ấy một cái cho bõ tức. Nếu không phải vì hệ thống công cụ tự động hồi phục công trình hư hại được tôi cài đặt từ năm ngoái đang âm thầm sửa chữa lại bản lề, thì có lẽ tôi đã làm thật rồi chứ chẳng đùa. Thay vào đó, tôi chỉ đành nén giận thở dài, buông một câu nhắc nhở cộc lốc: ``Có tay có chân đàng hoàng thì lần sau làm ơn mở cái cửa một cách bình thường giùm anh cái.''

Kaiho tặc lưỡi: ``Cứ đà này thì ngân sách tu sửa cơ sở vật chất của công ty cũng báo động đỏ theo \textit{holoX} mất thôi.''

Nàng Vịt nhìn thấy người bước vào như bắt được vàng, mừng rỡ reo lên: ``Aki-senpai! Cứu em với!''

Nàng Dị giới nghiêng đầu ngơ ngác: ``Có chuyện gì buồn phiền bực dọc về chị à em?''

Game bình phẩm: ``Dùng từ `phát điên' thì nghe có vẻ chuẩn xác hơn đấy. Nhưng yên tâm đi, đối tượng gây ức chế không phải là cô đâu, Rosenthal-san.''

Subaru gào lên: `\textbf{`Không hề nha! Cái tên ngốc lỳ lợm này cứ liên tục\dots}''

Nàng Tomboy hậm hực chỉ tay thẳng mặt Lui-chan - kẻ lúc này cũng vừa kịp nuốt trôi miếng bánh naan khổng lồ cuối cùng và đang ngồi chắp tay tỏ vẻ vô tội vạ - cố gắng giải thích đầu đuôi câu chuyện cho người tiền bối của mình nghe.

Bất thình lình, toàn bộ hệ thống đèn điện trong văn phòng vụt tắt ngấm. Một luồng ánh đèn sân khấu rực rỡ từ trên trần nhà rọi thẳng xuống vị trí của nàng Chim ưng, cùng lúc đó, một màn hình trình chiếu khổng lồ hiện lên lù lù ngay phía sau lưng em ấy, trên mặt màn hình chạy dòng chữ màu vàng chói lọi: ``Takane Lui - Khóa học quản lý tài sản chuyên sâu''. Còn tất cả chúng tôi đứng ở phía dưới bỗng chốc biến thành ``khán giả'', hay nói một cách chính xác hơn là đám ``sinh viên'' bất đắc dĩ.

\img{13-7i}

Nhìn cái khung cảnh này cứ y hệt như hồi năm ngoái tôi ngồi ngáp ngắn ngáp dài học môn Kinh tế kỹ thuật ở trường đại học vậy nhỉ? Cá nhân tôi thì chẳng lấy gì làm phiền lòng về chuyện này, cơ mà liếc sang bên cạnh thì có vẻ Subaru và Kaiho lại chẳng thoải mái dễ chịu chút nào. Riêng Aki thì\dots\ khỏi phải bàn cãi, cái gì lạ lẫm em ấy cũng háo hức muốn tìm hiểu, đúng chuẩn dạng ``chiếu mới chưa trải'' thứ thiệt.

Lui đứng trên bục giảng, chống tay đầy chuyên nghiệp: ``Mấy anh chị ngồi ở dưới đây đều đang có nhu cầu muốn đầu tư sinh lời đúng không ạ?''

Game ngơ ngác: ``\dots Cái quái gì đang diễn ra thế này?''

Nàng Vịt ôm trán: ``Ối dồi ôi. Lại giở chứng nữa rồi.''

Cô em gái song sinh khoanh tay khinh bỉ: ``Chị chả có nhu cầu nghe một kẻ đang báo động đỏ ngân sách dạy cách quản lý tài sản đâu.''

Tôi nhún vai cười nhẹ: ``Sao đâu? Nghe giảng giải trí cũng vui phết mà.''

Kaiho quay sang trừng mắt với tôi: ``Vui cái nỗi gì mà vui\dots\ Em cũng đã từng ngáp ngắn ngáp dài mấy bài triết kiểu này rồi đấy. Thà em đắp chăn đi ngủ hoặc là đi tranh luận Vật lý kỹ thuật còn sướng hơn.''

\dots Thì đúng là thế thật, bản thân tôi vốn dĩ là một kẻ cuồng học hỏi kiến thức mới mà. Với tôi, bất cứ thứ thông tin gì có thể thu nạp được vào đầu thì tôi đều sẵn sàng ``chơi'' tất tần tật. Trong cái ``khóa học'' tự phong đầy ngẫu hứng này, chính nàng Chim ưng kia sẽ đóng vai trò là giảng viên thuyết trình chính.

Lui dõng dạc: ``Bất kỳ ai cũng có thể thực hiện thành công việc quản lý tài sản cá nhân.'' Cái chân lý này thì hiển nhiên quá rồi, chẳng cần phải bàn cãi.

``Thế nhưng ngay bây giờ, để em trân trọng giới thiệu với các anh chị về một phương pháp tối ưu, cũng chính là một cơ hội đầu tư đỉnh cao đang được áp dụng rộng rãi trong khoảng thời gian gần đây.''

Aki hớn hở vỗ tay: ``Cái gì đây, cái gì đây? Hồi hộp quá đi mất!''

Game lo lắng lẩm bẩm: ``Tôi chỉ cầu mong nó không phải là một trò chơi chữ nhảm nhí nào nữa\dots''

Ngay lập tức, màn hình bài thuyết trình phía sau bật sáng, hiển thị hình ảnh rõ nét của một miếng nhựa ngăn cách đồ ăn trong hộp cơm trưa truyền thống (hay trong tiếng Nhật gọi là ``baran''), kèm theo một mũi tên màu đỏ chót chỉ thẳng tới biểu đồ tỷ giá hối đoái ngoại tệ.

\img{13-7j}

Vị Trợ lý ảo thở dài sườn sượt: ``Tôi biết ngay mà\dots\ Không trật đi đâu được\dots''

Kaiho đưa tay bóp trán, gân xanh lại nổi lên: ``Lại là chơi chữ kiểu ba-lần với baran. Mấy người thích tra tấn lỗ tai người khác bằng dad jokes lắm đúng không? Em thề là nó à dám nhả thêm một câu chơi chữ nhạt toẹt nào nữa thì tôi đấm nát cái bục giảng kia cho coi.''

Lui tươi cười hỏi, không hề nao núng gì trước sự hăm dọa của cô em gái song sinh: ``Mọi người ở đây chắc hẳn đều biết rõ cái miếng ngăn cách này là vật dụng gì rồi chứ ạ?''

Aki nhanh nhảu giơ tay: ``Là cái miếng nhựa chuyên dùng để trang trí và phân chia thức ăn trong hộp bento phải không em?''

Game bổ sung: ``Ý cô muốn nhắc tới mấy cái miếng lá nhựa màu xanh xanh có viền răng cưa để chia ngăn cơm với thức ăn đó á?''

Vị ``giảng viên'' gật đầu tán thưởng: ``Hoàn toàn chính xác! Hai người đúng là có vốn hiểu biết sâu rộng về đời sống đấy.''

Được giảng viên thuyết trình cất lời khen ngợi, nàng Dị giới tỏ vẻ ngượng ngùng bẽn lẽn, nở một nụ cười tươi tắn rạng rỡ như hoa hướng dương. Thế nhưng, vấn đề cốt lõi vẫn còn nằm chình ình ở đó: ``Nhưng mà rốt cuộc thì mấy cái miếng lá nhựa chia cơm hộp đó có liên quan quái gì đến nghệ thuật quản lý tài sản cơ chứ!?'' đó cũng chính xác là câu hỏi nhức nhối mà trợ lý ảo của tôi, Kaiho, cũng như toàn thể khán giả bất đắc dĩ chúng tôi đều đang muốn gào lên thắc mắc.

Và để giải đáp thắc mắc chung của cả phòng, Lui-sensei (tôi gọi trang trọng vậy cho phải phép) hắng giọng tiếp tục bài giảng của mình: ``Dạo gần đây, những hộp cơm trưa bento bày bán tại các cửa hàng tiện lợi Nhật Bản đang tạo nên một cơn sốt làm mưa làm gió trên toàn thế giới. Đặc biệt ở các quốc gia dầu mỏ giàu có, việc tiêu thụ món cơm này thậm chí đã vươn tầm trở thành biểu tượng khẳng định đẳng cấp của tầng lớp thượng lưu đấy ạ!''

\img{13-7k}

\dots Nghe cứ thấy sai sai và cấn cấn ở đâu đó trong câu chuyện này. Cả bốn chúng tôi đều lập tức cảm nhận được sự vô lý rành rành. Thế nhưng, nàng tóc vàng ngây thơ Aki vẫn hoàn toàn tin sái cổ vào những lời bịa đặt trơn tru đó, vỗ tay phành phạch với ánh mắt rực sáng đầy ngưỡng mộ.

Nàng DỊ giới reo lên: ``Thật luôn sao? Tuyệt vời quá đi mất!''

Subaru và Game đồng thanh hỏi vặn: ``Nghiêm túc đấy à?''

Tôi gãi đầu bất lực: ``Anh không nghĩ món cơm hộp bình dân này lại có sức ảnh hưởng vĩ mô đến thế đâu\dots''

Kaiho tặc lưỡi: ``Nghe mùi lừa đảo đa cấp nồng nặc thế này mà cũng có người tin cho được.''

Phớt lờ sự hoài nghi của đám đông, cô nàng tóc hồng-nâu tiếp tục thao thao bất tuyệt: ``Thêm vào đó, giới thượng lưu còn có thói quen sử dụng chính miếng nhựa baran như một tấm huân chương danh giá để chứng minh với thiên hạ rằng họ đã thưởng thức một hộp cơm thượng hạng như vậy.''

Để bài giảng thêm phần trực quan sinh động, Lui-sensei còn cẩn thận trưng ra một chiếc huy hiệu tự chế, với chiếc lá nhựa giả màu xanh được đính đằng sau một vật tròn nhỏ (nhìn kỹ thì trông y hệt cái nắp chai nhựa) ngay trên ngực áo. Tôi\dots\ vẫn hoàn toàn mù tịt, không tài nào liên kết nổi cái thứ đồ chơi thủ công này với bài toán kinh tế vĩ mô đang bàn.

Aki nghiêng đầu thích thú: ``Trông nó cũng có nét đáng yêu đấy chứ nhỉ\dots''

Lui lập tức tung hô: ``Đúng là senpai có gu thẩm mỹ tinh tế bậc nhất!''

Game phũ phàng bình phẩm: ``Còn tôi thì chỉ thấy nó nhảm nhí và rẻ tiền hết sức\dots''

Xuyên suốt từ đầu buổi thuyết trình đến giờ, chỉ có duy nhất Aki là học viên nhiệt huyết và hăng hái nhất, phần lớn cũng là vì được giảng viên liên tục tâng bốc khen ngợi. Còn lại, bốn chúng tôi đều đang ở trạng thái ``tâm hồn treo ngược cành cây'', những thông tin chúng tôi cố nhét vào đầu chỉ là một mớ bòng bong hỗn độn, càng nghe càng không hiểu em ấy này đang luyên thuyên cái quái gì.

Lui làm vẻ mặt tiếc nuối: ``Tuy nhiên, do các hộp cơm trưa có hạn sử dụng quá ngắn ngủi, nên những chiếc lá nhựa đi kèm thường được ví như những món trang sức xa xỉ đắt đỏ.''

Aki sụt sịt đồng cảm: ``Nghe buồn ghê á.''

Game cạn lời: ``\dots Thì người ta hoàn toàn có thể nhặt riêng cái lá nhựa ấy ra dùng lại được mà?''

Nàng Dị giới nghiêm túc phản bác: ``Nhưng làm như thế thì còn đâu giá trị của trang sức cao cấp nữa!''

Subaru quay sang huých tay tôi, thì thầm tuyệt vọng: ``\hl{Nè, từ nãy tới giờ anh có tiêu hóa nổi chữ nào em ấy nói không vậy?}''

Tôi thở dài sườn sượt: ``\hl{\dots Tin anh đi. Ngồi nghe giảng ba tiết Kinh tế - Chính trị liền tù tì ở trên giảng đường đại học có khi còn dễ thấm hơn vạn lần so với những lý thuyết quái đản mà em ấy đang nhồi nhét này.}''

Đúng lúc này, đôi mắt của thuyết giảng viên Chim ưng chợt lóe lên tia ranh mãnh: ``Đúng như Game-san vừa phân tích đấy. Nếu thứ đó đơn thuần chỉ là một miếng baran bình thường thì sao nhỉ?''

Nàng Dị giới Aki không ngần ngại giơ tay đáp lời, dù giọng điệu có phần hơi ngập ngừng: ``Thì\dots\ bất kỳ ai cũng có thể dễ dàng sở hữu được một cái ư?''

``Hoàn toàn chính xác!''

Tôi ngơ ngác chớp mắt: ``Ờm\dots\ Hả?''

Lui dõng dạc tuyên bố: ``Chính vì kẽ hở đó mà hiện tại, món phụ kiện baran cài áo - hay gọi nôm na là hàng baran chợ đen - đang trở thành kênh đầu tư siêu lợi nhuận hái ra tiền cực kỳ béo bở đấy ạ!''

Game nhướng mày phân tích: ``Nói tóm lại là chúng ta sẽ sản xuất hàng nhái bằng phương pháp thủ công dễ dàng hơn nhưng lại thu về biên độ lợi nhuận cao gấp bội, đúng chứ?''

``Chuẩn không cần chỉnh.''

Trên màn hình trình chiếu lúc này hiện ra sơ đồ minh họa một chiếc tàu chở hàng cỡ lớn chỉ thẳng đến đống lá nhựa giả, rồi nối tiếp đến chiếc huy hiệu; trong khi đó, nhánh nguồn cung chính thống đi từ hộp cơm bento đến baran rồi thành huy hiệu thì bị bôi xám xịt và gạch chéo đỏ lòm.

\img{13-7l}

Aki lo lắng hỏi: ``Nhưng nếu muốn nhập hàng kinh doanh như thế thì chắc sẽ phải cần nhiều vốn lắm, đúng không em?''

Lui mỉm cười bí hiểm: ``Đúng là như vậy. Một lô gồm một trăm chiếc như thế trên thị trường có giá lên tới một triệu yên lận!''

Nàng Dị giới giật mình: ``Đắt cắt cổ thế cơ á!''

Game gắt lên: ``Bán cái giá đó thì có khác gì ăn cướp giữa ban ngày đâu!''

Kaiho khoanh tay cười nhạt: ``Bán lá nhựa giá một vạn yên một chiếc, đúng là tư duy làm giàu không khó phiên bản lừa đảo.''

Ánh đèn văn phòng bật sáng trở lại bình thường, nhưng Lui và Aki vẫn say sưa tiếp tục màn vấn đáp mua bán đầy kịch tính.

Nàng Chim ưng hạ giọng dụ dỗ: ``Nhưng mà chị đừng có lo lắng quá. Vì chúng ta đều là người một nhà ở hololive, em sẽ ưu tiên bán lẻ số baran đó cho riêng chị với mức giá cực kỳ phải chăng!''

Aki sáng rực mắt: ``Thật vậy sao!?''

Tôi lắc đầu ngán ngẩm: ``Nghe y chang mấy cái khóa học làm giàu lùa gà trên mạng nhỉ\dots''

Lui ghé sát vào tai đối phương, thì thầm đầy mờ ám: ``\hl{Ưu đãi độc quyền này chỉ dành riêng cho một mình chị thôi đó, senpai à.}''

Ái chà chà, lại còn bày đặt giấu giấu diếm diếm tạo cảm giác khan hiếm cơ đấy! Mà thực lòng nếu có mời chào miễn phí chắc tôi cũng chẳng thèm đụng vào ba cái lá nhựa chia cơm hộp ấy làm gì.

Nàng Dị giới đứng đăm chiêu suy nghĩ một lúc lâu, rồi đột ngột đưa ra một quyết định vô cùng táo bạo: ``Chốt đơn! Chắc chị sẽ mua mở hàng mười cái.''

Nàng Chim ưng lập tức chắp tay cảm kích: ``Vậy tổng thiệt hại nhẹ nhàng của chị sẽ là 100 nghìn yên ạ!''

Trợ lý ảo của tôi đổ mồ hôi hột: ``Dễ dãi chốt kèo thế cũng được luôn nữa hả trời?'' còn Kaiho thì chỉ biết khoanh tay tặc lưỡi: ``Còn lạ gì cái kiểu con nai tơ ngơ ngác từ Aki-chan được chứ.''

Chưa dừng lại ở đó, Lui tiếp tục nắm chặt lấy đôi bàn tay của người tiền bối, dồn dập dẫn dụ Aki bằng những lời đường mật có cánh: ``Chị không chỉ đơn thuần bỏ tiền ra mua baran đâu\dots'' câu nói lấp lửng khiến đối phương ngạc nhiên.

``Chị còn đang trực tiếp đầu tư vào một tương lai chắc chắn sẽ thành công rực rỡ đấy ạ!'' Lui-sensei tuyên bố hùng hồn, xung quanh hiệu ứng hàng chục bông hoa nở rộ lấp lánh như mở ra một viễn cảnh xán lạn đang vẫy gọi ở phía trước.

\img{13-7m}

Nàng tóc vàng đến nước này thì đã thực sự tin đến sái cổ, thậm chí trong đôi mắt em ấy còn hiện rõ biểu tượng đồng tiền lấp lánh: ``Chị đồng ý! Chị sẽ tham gia dự án này!''

Đến giây phút này thì tôi đã có thể đưa ra kết luận cuối cùng: đây đích thực là một khóa học lừa đảo đa cấp biến tướng chứ chẳng phải khóa học quản lý tài sản hay đầu tư kinh tế gì sất! Subaru và hai anh em tôi đứng ngoài lo lắng cho sự ngây thơ của Aki đến toát cả mồ hôi hột. Tôi nghĩ có lẽ Lui cuối cùng cũng đã lộ rõ bản chất thật: em ấy này ranh ma và xảo quyệt hơn tôi tưởng tượng rất nhiều. Hoặc có lẽ là do Aki quá đỗi ngây thơ, hoặc là cả hai yếu tố cộng hưởng lại.

Subaru hoảng loạn kêu lên: ``Chết dở rồi, chị ấy tin sái cổ luôn rồi kìa!''

Kaiho thở dài ngán ngẩm: ``Rơi thẳng vào bẫy lùa gà luôn rồi chứ còn gì nữa\dots''

Aki chắp tay thành kính: ``Nếu những gì mà vị giáo sư tài ba này phân tích là đúng, thì chắc chắn tương lai sẽ giàu to!''

Nàng Vịt đổ mồ hôi hột đầm đìa: ``Chị ấy bị tẩy não say mê luôn rồi!''

Game thở dài: ``Cô thực sự khờ dại đến mức khó tin đấy, Rosenthal ạ\dots''

Subaru gào lên: ``Với lại, ở cái văn phòng này làm quái gì có ông giáo sư nào đâu cơ chứ!''

Nhận thấy nàng Vịt vẫn còn giữ được sự tỉnh táo nhất định, nàng Chim ưng liền chuyển mục tiêu, tiếp tục giở giọng ngon ngọt để tiếp cận em ấy: ``Chị có muốn cùng tham gia đầu tư sinh lời với tụi em không, Subaru-senpai?''

Lui cứ thế từ từ, chầm chậm từng bước tiến sát lại gần. Chiếc bóng đen ngòm của em ấy đổ dài rọi lên người nàng Vịt đang run rẩy sợ sệt cứ mỗi lúc một lớn dần, tỷ lệ thuận với nỗi khiếp đảm tột độ đang dâng trào trong lòng Subaru.

Đến giới hạn chịu đựng, khuôn mặt Subaru tái mét cắt không còn giọt máu, cô nàng vùng lên phản kháng mạnh mẽ: ``Ở ngoài xã hội kia vẫn còn đầy rẫy những cạm bẫy xấu xa đáng sợ lắm!''

Nói đoạn, cô một tay túm chặt kéo lấy cô nàng tóc vàng vẫn còn đang mải mê đắm chìm trong giấc mơ làm giàu triệu phú, ba chân bốn cẳng cắm cổ chạy thục mạng ra khỏi văn phòng, vừa chạy vừa hét toáng lên thất thanh: ``\textbf{Cứu với! Cái Hội kín Bí mật này đáng sợ và nguy hiểm quá đi mất!!}''

\img{13-7n}

Trong căn phòng rộng lớn lúc này chỉ còn lại tôi, Kaiho và Lui. Ngay khoảnh khắc bóng dáng của hai người kia vừa khuất sau cánh cửa, nàng tóc nâu liền quay ngoắt về phía tôi và\dots\ thè lưỡi tinh nghịch, nháy mắt: ``È-hé\~{}''

Tôi thở dài ngao ngán: ``Thôi xin em, đừng có mà `è-hé' vô tội vạ trước mặt anh như thế.''

Trợ lý của tôi cũng lên tiếng bày tỏ thái độ không đồng tình với cách hành xử dọa người của em ấy: ``Cô vừa mới làm cho họ sợ chết khiếp suýt đứng tim đấy, Takane-san ạ.''

Kaiho khoanh tay lắc đầu: ``Diễn xuất nhập tâm cỡ này thì đi làm nhân viên kinh doanh đa cấp chắc hốt bạc.''

Tất nhiên, đáp lại những lời nhắc nhở đó chỉ là một khuôn mặt hồn nhiên ngây thơ vô số tội như em bé của Lui-chan: ``Em có cố ý dọa ai đâu thật mà!''

\dots Đây đúng là một món quà tinh thần đầy bất ngờ nhân dịp kỷ niệm tuổi hai sáu lăm của tôi. Thật không thể ngờ cô nàng quản lý của holoX lại sở hữu tài năng thao túng tâm lý đỉnh cao đối với những ``chú nai vàng ngơ ngác'' đến mức độ này. Quả đúng là một nhân vật nguy hiểm không thể xem thường!

%==========================================TRUYỆN PHỤ=========================================
\omk

Bầu không khí bên trong văn phòng bỗng chốc trở nên tĩnh lặng đến kỳ lạ ngay sau khi Subaru kéo cổ Aki cắm đầu chạy biến ra ngoài hành lang. Khói bụi tấu hài tan biến, để lại tôi, trợ lý Game và Lui đứng đối diện nhau giữa gian phòng trống trải. Về phần Kaiho, vì còn phải tất bật kiểm tra nốt lịch trình buổi tối cho hội \hls\ nên em ấy đã vội vã rời sang khu vực căn mancave từ lúc màn rượt đuổi vừa kết thúc, hẹn gặp lại tôi ở ga tàu để cùng về nhà mừng sinh nhật đôi của hai anh em.

Lui-chan thu lại vẻ mặt vô tội, nghiêng đầu nhìn tôi đầy tò mò: ``Mà công nhận cũng lạ thật ha, Kuon-senpai?''

Tôi nhướng mày hỏi lại: ``Lạ ở chỗ nào cơ?''

Lui khoanh tay, ánh mắt dò xét: ``Trong suốt cả quá trình diễn kịch vừa rồi, anh có vẻ như hoàn toàn chẳng hề có chút cảm xúc sợ hãi hay nao núng gì trước những mánh khóe thao túng của em cả?''

Nàng Chim ưng nhìn chằm chằm vào tôi với hàng loạt dấu chấm hỏi to đùng lơ lửng trên đầu. Quả thực, xuyên suốt quá trình mà em ấy giở chiêu ``lùa gà'' cô nàng ngây thơ Aki, nếu Subaru sợ đến tái mét mặt mày, Kaiho thì ngán ngẩm buông những lời cay nghiến, còn Game thì cạn lời tặc lưỡi, thì chỉ có duy nhất bản thân tôi là vẫn giữ được sự bình thản đến mức lạnh lùng.

Tôi nhún vai, điềm nhiên đáp lại bằng một giọng điệu vô cùng thực tế: ``Bởi vì ở cái nơi xuất thân của anh, chiêu trò lừa đảo đa cấp biến tướng kiểu này nhan nhản ngoài đường nhiều như nấm mọc sau mưa ấy mà. So với bọn họ, trình độ của em vẫn còn thuộc dạng nghiệp dư lắm, chưa đủ tầm để hù dọa được anh đâu.''

Lui bật cười trừ, đưa tay gãi nhẹ má: ``Hì hì\dots\ Bị anh bắt bài sạch sành sanh rồi à?''

Tôi gật gù nhận xét khách quan: ``Nhưng mà công nhận diễn xuất nhập tâm phết đấy, trông em cứ như thể một doanh nhân thành đạt thực thụ đi thuyết giảng vậy. Cũng khá lắm.''

Game lên tiếng từ trong chiếc đồng hồ đeo tay, vặn vẹo thực tế: ``Với cả, kể cả nếu Rosenthal-san có thực sự rút ví đưa tiền thật, thì cô tính xoay xở sản xuất hàng loạt kiểu gì với mấy cái lá nhựa chia cơm hộp đó?''

Lui xua tay cười xòa: ``Tất nhiên là làm sao mà triển khai thật được chứ ạ! Cơ bản là em cũng chẳng ngờ nổi là chị ấy lại ngây thơ tin sái cổ và sập bẫy nhanh đến thế cơ đấy!''

Tôi thở dài sườn sượt: ``Em ấy vốn dĩ là một tờ giấy trắng tinh khôi mà. Tâm hồn còn ngây thơ trong sáng lắm, sau này chắc chắn còn phải va vấp học hỏi trường đời dài dài.''

Vị trợ lý ảo chép miệng châm chọc: ``Nhân tiện thì, cô chơi chữ nhạt nhẽo thật đấy, Takane-san. Trên đời thiếu gì thứ không chọn lại đi chọn đúng cái miếng nhựa baran để ví von tài chính.''

Lui khẽ cười gượng gạo, cúi đầu nhận lỗi: ``Xin lỗi ạ\dots\ Thói quen nhiễm mấy trò đùa nhạt của các ông bố trung niên rồi.''

Tôi vỗ tay một cái bốp để kéo mọi người quay trở lại quỹ đạo chính: ``Thôi, dẹp chuyện đùa cợt sang một bên đi. Giờ thì\dots\ chúng ta bắt đầu tìm cách giải quyết tiếp vấn đề cốt lõi nhé?''

Game hỏi lại cho chắc: ``Ý cậu là cái bài toán nan giải về việc tìm thêm nguồn thu nhập cho Hội kín X á?''

``Đúng, chính xác là cái mớ bòng bong đó đấy.''

Lui chớp mắt: ``À, giờ em mới sực nhớ ra mục đích ban đầu của buổi gặp hôm nay\dots''

Thế là cả ba người chúng tôi lại tiếp tục chụm đầu vào bàn bạc, vắt óc suy nghĩ về hàng loạt phương án để cân đối lại bảng thu chi bết bát của tổ chức \textit{holoX}. Chúng tôi tranh luận, gạch xóa và tính toán suốt gần nửa tiếng đồng hồ ròng rã, thế nhưng cuối cùng vẫn chẳng thể nào thống nhất được một hướng đi khả thi nào.

Nàng Chim ưng gục đầu thở dài ngao ngán: ``Haiz\dots\ Xem ra phương án tăng thu nhập không ổn rồi nhỉ.''

Game phân tích ngắn gọn: ``Căn cứ theo thực trạng nhân sự hiện tại của tổ chức cô, tôi khẳng định là không còn phương án tạo doanh thu nào khả dĩ hơn nữa đâu.''

Lui suy tư một giây, rồi khẽ mỉm cười: ``Vậy chắc là chúng ta buộc phải kích hoạt Kế hoạch B thôi nhỉ?''

Tôi tò mò hỏi: ``Kế hoạch B là cái quái gì nữa vậy?''

Một lần nữa, đôi mắt chim ưng màu vàng rực rỡ với biểu tượng chữ thập đỏ hồng sắc lẹm lại hiện lên trong đáy mắt em ấy, đi kèm với một thần thái nghiêm túc đến đáng sợ: ``Thắt chặt chi tiêu tuyệt đối!''

Cả Game và tôi đều đồng loạt gật gù tán thành: ``Cách này nghe có vẻ thực tế và khả thi hơn vạn lần đấy.''

Vị trợ lý ảo hỏi tiếp: ``Vậy cụ thể thì cô dự định sẽ cắt giảm ngân sách cho những hạng mục nào đầu tiên?''

Lui đáp gọn lỏn: ``Tất cả những thứ vớ vẩn thực sự không cần thiết.''

Tôi búng tay cái tách: ``OK, ý tưởng rất hay. Em thử liệt kê chi tiết cụ thể hơn xem nào?''

Nàng Chim ưng đưa ngón tay chỉ thẳng vào những phần chiếm diện tích lớn trên chiếc biểu đồ tròn mà chúng tôi đã xem lúc đầu: ``Khoản gà rán công nghiệp, chi phí nhập cà tím, phí tiêu vặt linh tinh\dots\ Nói tóm lại là dẹp sạch toàn bộ mấy thứ vô bổ mà Sếp tổng của bọn em hay vung tay quá trán chi tiêu!''

Tôi vô thức đổ mồ hôi hột đầm đìa: ``Thật sự là thói quen tiêu xài của người đứng đầu tổ chức các em đang có vấn đề nghiêm trọng đấy\dots''

Game lạnh lùng bồi thêm một câu: ``Đến cả cô em gái út hay mộng mơ của Kuon-user cũng chẳng bao giờ tiêu xài hoang phí tùy tiện đến mức độ vô lý như thế này.''

Lui-chan bật cười khúc khích, buông một câu đùa bâng quơ: ``Vậy chắc là bà chủ nhà bọn em phải cắp sách sang học hỏi kinh nghiệm quản lý từ em gái của anh rồi nhỉ?''

Thế nhưng, đáp lại sự hóm hỉnh đó chỉ là nụ cười gượng gạo sượng trân đến cứng cả cơ mặt của tôi: ``Ừ-Ừm? Chắc là vậy rồi đấy\dots''

Nhận thấy sự im lặng ngượng ngùng bao trùm cả căn phòng, em ấy lập tức gục mặt xuống bàn, cất giọng thất vọng não nề: ``Em xin lỗi anh\dots\ Câu đùa vừa rồi thực sự không buồn cười chút nào nhỉ?''

Tôi hoảng hốt xua tay thanh minh lia lịa: ``K-Không hề có chuyện đó đâu! N-Nó cũng hài hước và duyên dáng ra phết đấy chứ!''

Dù thú thực trong lòng thì\dots\ câu đùa vừa rồi còn nhạt hơn cả nước ốc luộc. Bản thân tôi vốn dĩ là một người rất dễ tính và dễ cười, thế nhưng chẳng hiểu sao đứng trước mấy màn chơi chữ của em ấy này, dây thần kinh hài hước của tôi lại hoàn toàn đình công không thể hoạt động nổi.

\ct

Đồng hồ điểm giờ tan tầm, đã đến lúc bọn tôi thu dọn đồ đạc ra về. Lui quay trở lại hang ổ bí mật của Hội kín nằm sâu bên trong tòa nhà công ty, còn tôi thì như mọi khi, vác ba lô ra ga tàu để cùng cô em gái song sinh Kaiho trở về nhà.

Lui vẫy tay tiễn tôi ra cửa: ``Anh về nhà cẩn thận nha, Kuon-senpai.''

Tôi gật đầu đáp lại: ``Ừ, em cũng nghỉ ngơi sớm đi. Hẹn ngày mai gặp lại.''

Tôi đưa tay mở cánh cửa phòng vừa được hệ thống tự động sửa chữa lại nguyên vẹn như mới và chuẩn bị sải bước ra hành lang, thế nhưng ngay khoảnh khắc đó, một tiếng gọi bất ngờ vang lên từ phía sau: ``Khoan đã, Kuon-senpai ơi!''

Tôi khựng lại, ngoảnh đầu nhìn về phía em ấy: ``Có chuyện gì thế?''

Nàng Chim ưng từ tốn sải bước tiến tới gần, hai tay nâng niu đưa ra trước mặt tôi một chiếc bánh cốc (cupcake) sô-cô-la được trang trí khá tỉ mỉ và đẹp mắt. Ờm\dots\ tự dưng lại tặng bánh ngọt vào giờ này là có ý đồ gì đây?

Và rồi, điều bất ngờ mà tôi nằm mơ cũng không thể lường trước được đã xảy đến.

Lui nở một nụ cười rạng rỡ, dõng dạc cất lời: ``Chúc mừng sinh nhật lần thứ hai mươi sáu của anh, Kuon-senpai! À, cả lời chúc mừng gửi tới Kaiho-senpai nữa nhé!''

Tôi trố mắt kinh ngạc đến mức đứng hình mất vài giây, miệng há hốc không khép lại nổi: ``S-Sao\dots\ làm thế nào mà em lại biết được hôm nay là sinh nhật của anh em bọn anh vậy!?''

Lui nháy mắt tinh nghịch: ``Em là người quản lý tổng hợp kiêm phụ trách tình báo của \textit{holoX} mà, anh quên rồi sao? Với cả, chính Sếp tổng của bọn em đã ra lệnh bắt buộc em phải điều tra và nắm rõ mọi thông tin hồ sơ về anh đấy.''

Tôi gãi đầu bất lực, nheo mắt gặng hỏi: ``Nhưng mà dù em có điều tra hồ sơ nhân sự công khai của công ty thì ngày sinh của Kaiho cũng thuộc diện bảo mật riêng của nhánh \hls\ kia mà, làm sao em có thể nắm tường tận đến từng chân tơ kẽ tóc như thế được?''

Lui nghiêng đầu, buột miệng đáp một câu tỉnh bơ: ``À\dots\ thì cũng nhờ thông tin mạng lưới của nhóm thu thập được thôi ạ. Mà chủ yếu là từ nguồn tin của bộ ba nhà Chó cung cấp đấy chứ\dots''

Tôi giật mình khựng lại, mắt mở to kinh ngạc: ``Hả\dots? Bộ ba nhà Chó nào cơ?''

Dường như nhận ra mình vừa lỡ lời để lộ một thông tin tuyệt mật không nên nói, nàng Chim ưng giật bắn mình, vội vàng lấy hai tay bịt chặt miệng lại. Đôi mắt em ấy đảo liên hồi tỏ vẻ lúng túng tột độ, vội vàng xua tay chữa cháy trong hoảng hốt: ``A\dots\ Á! K-Không có gì đâu ạ! Em\dots\ em chỉ đang nói đùa theo mấy mật danh trinh sát nội bộ của bọn em thôi mà! Anh đừng bận tâm làm gì nha, haha\dots''

Tôi và Game đưa mắt nhìn nhau đầy hoài nghi.

Trong tâm trí tôi lúc này, cụm từ ``nhà Chó'' hay ``ba mẹ con nhà Chó'' chỉ có thể gợi nhớ về một ký ức đượm buồn ở xưởng cơ khí của gia đình Hikawa\footnote{Xem thêm diễn biến trong {\flaregothic HISTOIRE \freshbot 40}, thuộc quyển Ký Túc Xá.}: đó là chú chó mẹ Mi trung thành đã gắn bó với gia đình tôi từ thuở nhỏ, cùng hai đứa con của nó là Tô và Ki. Thế nhưng\dots\ Mi đã già yếu và qua đời vào một buổi chiều tháng trước dưới gốc cây quen thuộc trong vườn xưởng, còn Tô và Ki thì sau đó không lâu cũng đột ngột mất tích một cách bí ẩn, khiến Suzuno khóc ròng suốt mấy ngày liền và cả nhà cất công đi tìm khắp các con phố ở Kawachi mà chẳng thấy tăm hơi.

Mi đã yên nghỉ nơi vĩnh hằng, còn Tô và Ki có lẽ đã lưu lạc nơi nào đó hoặc bị kẻ xấu bắt trộm mất.

Chính vì thế, tôi chỉ tặc lưỡi tự nhủ chắc em ấy này lại đang học đòi nói mấy thuật ngữ ám hiệu kỳ quặc của các tổ chức bí mật trên phim ảnh.

Trợ lý Game cất giọng nhắc nhở từ cổ tay: ``Chắc lại là mấy trò đặt mật danh hoang tưởng của cái Hội kín này thôi. Kuon-user đừng tốn tế bào thần kinh để suy luận làm gì cho mệt óc.''

Thấy tôi không gặng hỏi thêm, Lui thở phào nhẹ nhõm một hơi dài như trút được gánh nặng, vội vàng chuyển chủ đề: ``Với cả, nhân tiện đây cho em gửi lời cảm ơn chân thành tới anh vì đã cất công tắm rửa sạch sẽ cho Chloe nhé.''

Tôi ngơ ngác chớp mắt: ``Ờm\dots? Vụ đó thì có gì đâu\dots''

Lui bật cười: ``Em không rõ anh đã sử dụng biện pháp mạnh tay nào, nhưng mà tuần vừa rồi gặp lại em cảm thấy cậu ta thơm tho sạch sẽ đến mức bất ngờ luôn ấy. Gặng hỏi mãi mới biết hóa ra là do một tay anh ra tay xử lý.''

Tôi cười khổ, gãi má ngượng ngùng: ``\dots Ờm, anh\dots\ cảm ơn lời khen của em? Mặc dù anh nghĩ sau cái vụ cưỡng chế tắm rửa hôm đó thì em ấy cũng đâm ra ghét cay ghét đắng anh rồi\dots''

Lui khẽ cười đầy ẩn ý: ``Em nghĩ chắc hai người sẽ sớm làm lành lại với nhau thôi mà. Thôi, muộn rồi, anh mau đi về kẻo Kaiho-senpai đợi lâu nha!''

Tôi gật đầu nhận lấy chiếc bánh: ``Er\dots\ Rồi, cảm ơn em nhiều nhé.''

Tôi hối hả sải bước rời khỏi văn phòng, mồ hôi hột túa ra đầm đìa trên trán. Đáng sợ thật sự, bây giờ danh tính của một thằng quản lý quèn như tôi thậm chí còn lọt vào tầm ngắm điều tra của cả một tổ chức bí mật cơ đấy.

Game nghiêm túc nhắc nhở: ``Tôi nghĩ từ nay về sau chúng ta phải thực sự đề cao cảnh giác với cô nàng Takane này.''

Tôi thở dài đồng tình: ``\dots Ý kiến hoàn toàn xác đáng. Chúng ta chẳng thể nào lường trước được em ấy còn đang ấp ủ âm mưu gì đối với anh em nhà mình đâu--''

*TING!*

Bất chợt, tiếng chuông thông báo tin nhắn trên điện thoại vang lên lanh lảnh giữa hành lang vắng. Tôi khựng lại, rút chiếc điện thoại từ trong túi quần ra mở khóa màn hình.

Đó là một tin nhắn lạ vừa được gửi đến. Người gửi hiển thị một cái tên ngắn gọn nhưng đầy sặc mùi nguy hiểm: La+.

Tôi nhíu mày bấm mở nội dung tin nhắn ra đọc. Bên trong dòng văn bản được viết với một giọng điệu vô cùng ngạo nghễ và hống hách:

\txtbx{
    ``Oi, Hikawa Kuon!\\
    Có vẻ như ngươi cũng đã lần lượt chạm trán với bốn thành viên chủ chốt trong Hội kín của bọn ta rồi nhỉ? Đám thuộc hạ của ta dạo gần đây hay kể lể về ngươi nhiều lắm đấy.\\
    Iroha thì lúc nào cũng khép nép sợ sệt uy quyền của ngươi, Chloe thì vừa ghét cay ghét đắng lại vừa dè chừng ngươi, rồi đến cả Koyori cũng đâm ra hâm mộ ngươi ra mặt nữa! Không biết tên XO kia hôm nay sẽ báo cáo lại những gì cho ta đây!\\
    Xem chừng ngươi cũng là một đối thủ đáng gờm ra phết đấy, dám cả gan làm xáo trộn cả trật tự của Hội kín bọn ta nữa kìa. Ta dám chắc là ngươi đã sử dụng mánh khóe mờ ám gì đó rồi\\
    Dù sao đi nữa, tuần tới đây đích thân ta sẽ trực tiếp xuất đầu lộ diện gặp mặt ngươi, để tận mắt kiểm chứng xem Hikawa Kuon nhà ngươi thực sự có bản lĩnh gì. Trăm nghe không bằng một thấy mà, nhỉ?*\\
    Ta sẽ hẹn gặp mặt quyết đấu với ngươi vào ngày hôm đó! Đừng có nghĩ đến chuyện hèn nhát chạy trốn! Ta nắm rõ hết tất cả mọi bí mật về ngươi đấy!\\
    Ta thực sự, thực sự rất mong chờ được đối mặt trực diện với ngươi vào ngày đó lắm đấy!\\
    \flushright La+ - Thủ lĩnh tối cao của Hội kín holoX''
}

\dots Đọc xong từng dòng chữ sặc mùi thách thức kia, tự dưng tôi có cảm giác như thể mình đang chuẩn bị vác kiếm đi đánh trùm cuối trong một tựa game nhập vai vậy. Có lẽ nào ngay từ đầu, em ấy đã cố tình cử bốn thành viên trong Hội lần lượt xuất hiện trước mặt tôi là để âm thầm do thám tình hình thực lực của tôi hay sao?

Rốt cuộc thì cô nhóc này đang cần gì ở một thằng quản lý bình thường như tôi? Và tại sao em ấy lại khao khát muốn gặp mặt tôi đến mức như vậy?

Trong lúc tôi còn đang đứng ngẩn người suy nghĩ mông lung về mọi thứ xoay quanh cái Hội kín kỳ bí kia, trợ lý ảo của tôi bỗng cất giọng gọi dồn dập: ``Này, Kuon-user, Kuon-user!''

Tôi giật mình thoát khỏi dòng suy nghĩ:
``\dots H-Hả? Có chuyện gì thế?''

Game hất giọng nhắc nhở: ``Có vẻ như cô nhóc đó vừa mới gửi tiếp một tin nhắn bổ sung điều gì nữa kìa.''

Tôi vội vàng cúi xuống nhìn lại màn hình điện thoại. Quả nhiên, tài khoản mang tên La+ vừa mới gửi thêm cho tôi một đoạn tái bút ngắn gọn:

\txtbx{
    ``Tái bút: Mừng ngươi bước sang tuổi hai mươi sáu vui vẻ! Hy vọng ngươi sẽ thích chiếc bánh cốc sô-cô-la do chính tay ta làm! Đây là lần đầu tiên trong đời ta được tên Lui kia hướng dẫn làm bánh đấy nhé!\\
    À, tiện thể ta cũng đã tự ý thêm tài khoản của ngươi vào nhóm chat nội bộ của Hội rồi đấy! Nếu lúc nào rảnh rỗi thì cứ vào trò chuyện cùng bọn ta, ha?''}

\dots Đọc đến những dòng chữ ngây ngô đáng yêu này, tự dưng suy nghĩ trong đầu tôi về em ấy lại xoay chuyển một góc một trăm tám mươi độ.

Tôi có cảm tưởng\dots\ thay vì muốn tìm cách tiêu diệt kẻ thù, thì thực chất cô bé này đang ngưỡng mộ và muốn làm thân với tôi thì đúng hơn. Cơ mà ngặt một nỗi, bản thân tôi thì có cái điểm gì nổi bật để khiến cho người ta phải xiêu lòng cơ chứ? Tôi suy cho cùng cũng chỉ là một thằng nhân viên văn phòng quèn cắm đầu đi làm bảy ngày một tuần để kiếm cơm manh áo nuôi gia đình thôi mà!

``Anh đứng ngẩn ra đấy làm gì mãi thế, Aniki? Tàu sắp tới trạm rồi kìa!'' tiếng gọi sắc lạnh quen thuộc của Kaiho vang lên từ phía cổng chính tòa nhà kéo tôi trở về với thực tại.

Tôi vội vàng cất chiếc điện thoại vào túi, nâng niu hộp bánh cốc sô-cô-la trên tay rồi rảo bước chạy về phía cô em gái song sinh đang đứng đợi, khẽ cười thầm trong bụng.

Chậc\dots\ Thật sự là tôi chẳng tài nào đoán nổi trong đầu Laplus-chan đang suy nghĩ những gì nữa. Mà thôi kệ đi, cứ để đến khi nào ngày gặp mặt đó thực sự tới thì hẵng hay.

Sinh nhật tuổi hai mươi sáu của anh em chúng tôi có khác. Thật đúng là một ngày kỷ niệm ngập tràn những tình huống bẻ lái oái oăm khó đỡ!

\end{document}